\documentclass[conference]{IEEEtran}
\IEEEoverridecommandlockouts

% ---- standard IEEE packages ----
\usepackage{cite}
\usepackage{amsmath,amssymb,amsfonts}
\usepackage{algorithmic}
\usepackage{graphicx}
\usepackage{textcomp}
\usepackage{xcolor}
\usepackage{booktabs}
\usepackage{hyperref}
\usepackage{url}
\def\BibTeX{{\rm B\kern-.05em{\sc i\kern-.025em b}\kern-.08em
    T\kern-.1667em\lower.7ex\hbox{E}\kern-.125emX}}

\begin{document}

\title{Unsupervised Neural Networks for Multi-Genre \\
        Symbolic Music Generation: \\
        An LSTM / VAE / Transformer / RLHF Study}

\author{
\IEEEauthorblockN{Aneela Musarrat \\ID: 22101385\\Md Mamun Molla\\ID: 22201975\\CSE425: Neural Networks}
\IEEEauthorblockA{\textit{Department of Computer Science and Engineering} \\
\textit{24 April 2026} \\
Dhaka, Bangladesh}
}

\maketitle

\begin{abstract}
We investigate four unsupervised deep-learning approaches for symbolic
multi-genre music generation on a common vocabulary of 16\textsuperscript{th}-note
monophonic MIDI tokens. The four tasks form a pedagogical progression:
(i)~an LSTM-based sequence autoencoder trained on a single genre,
(ii)~a $\beta$-variational autoencoder with genre conditioning, 
(iii)~an autoregressive Transformer decoder, and
(iv)~a reinforcement-learning-from-human-feedback (RLHF) stage that
fine-tunes the Transformer policy against a learned preference reward.
Every model is implemented in pure NumPy+Autograd with analytic gradients,
trained on a five-genre corpus of 100~MIDI pieces (282 training windows), and
compared against a random-note and first-order Markov baseline using four
evaluation metrics: reconstruction loss, perplexity, rhythm diversity and
inter-sample pitch-histogram distance. The Transformer attains a final
per-token perplexity of~5.96, and RLHF further raises the combined musicality
reward from 0.56 to 0.67. All source code, 48~generated MIDI samples, training
plots and a 10-listener survey CSV accompany the report.
\end{abstract}

\begin{IEEEkeywords}
music generation, LSTM, variational autoencoder, Transformer, RLHF,
unsupervised learning, symbolic music, MIDI
\end{IEEEkeywords}

% ===========================================================================
\section{Introduction}
% ===========================================================================

Music is a structured temporal signal whose patterns --- melody, harmony,
rhythm and genre-specific stylistic distributions --- emerge from statistical
regularities in large corpora. Labelled music data is expensive, motivating
an \emph{unsupervised} approach that learns musical structure from raw
event sequences. Given a MIDI piece as a sequence of symbolic events
$X = \{x_1, x_2, \dots, x_T\}$, we target the generative distribution
\begin{equation}
p_\theta(X) \,=\, \prod_{t=1}^{T} p_\theta(x_t \mid x_{<t})
\label{eq:ar}
\end{equation}
so that novel pieces can be drawn as $\hat{X} \sim p_\theta(X)$.

This work implements four complementary models that capture
progressively richer aspects of $p_\theta(X)$: a recurrent autoencoder
(a deterministic compression view), a variational autoencoder (a probabilistic
generative view with a continuous latent space), a Transformer decoder
(explicit autoregressive factorisation with long-range attention), and
REINFORCE-style RLHF on top of the Transformer (preference alignment).

\paragraph*{Contributions}
\begin{enumerate}
  \item A self-contained NumPy+Autograd implementation of all four models
        with analytic gradients, suitable for environments without GPU
        support.
  \item A unified symbolic-music tokeniser that exports bidirectionally
        between MIDI files and integer token streams.
  \item End-to-end reproducibility: one command regenerates the corpus,
        trains every task, writes 48~MIDI samples, all training plots and
        the final metric table.
  \item A transparent hybrid reward model combining three rule-based
        musicality features with a linear regressor fitted on a
        10-listener preference survey, enabling reproducible RLHF on a
        course-project scale.
\end{enumerate}

% ===========================================================================
\section{Data Preparation}
% ===========================================================================

\subsection{Dataset}
The project is designed for any of the MAESTRO, Lakh MIDI or Groove MIDI
datasets~\cite{maestro,lakh,groove}; because those archives are several
gigabytes, we also ship a small \emph{synthetic} corpus generator
(\texttt{src/preprocessing/synthetic\_corpus.py}) that produces 20
genre-faithful MIDI files per genre using hand-crafted scale, tempo and
rhythm templates. The total corpus therefore consists of
\textbf{100 MIDI files} across five genres (classical, jazz, rock, pop,
electronic).

\subsection{Tokenisation pipeline}
Each \texttt{.mid} file is loaded with \texttt{pretty\_midi}, normalised to a
16-steps-per-bar grid and reduced to a monophonic stream by retaining the
highest-velocity pitch active at every step. The resulting integer tokens
cover an 88-pitch piano range, plus a \texttt{PAD}~(0) and \texttt{REST}~(1)
symbol, giving a vocabulary of up to 90 entries. The stream is sliced into
overlapping windows of length $T\!=\!64$ (four bars), yielding
\textbf{282 training windows}.

% ===========================================================================
\section{Task 1 --- LSTM Autoencoder}
% ===========================================================================

\subsection{Architecture}
The encoder is a single-layer LSTM that consumes an embedding $e_t \in
\mathbb{R}^{32}$ and produces a final hidden state $h_T \in \mathbb{R}^{64}$
that is projected to a latent code $z \in \mathbb{R}^{32}$:
\[
z = W_z \, h_T + b_z.
\]
The decoder LSTM receives the concatenation $[e_{t-1}; z]$ at every
time-step and a linear head maps its hidden state to logits over the
vocabulary.

\subsection{Loss}
With true label $x_t$ we minimise the per-step negative log-likelihood
\begin{equation}
\mathcal{L}_{AE}=-\frac{1}{T}\sum_{t=1}^{T}\log p_\theta(x_t\mid z,x_{<t}).
\end{equation}

\subsection{Training setup}
Classical-genre slice only (56~windows), 25~epochs, batch~16, Adam with
learning rate $5\!\times\!10^{-3}$. Training finished in 16.6~s on CPU.

\subsection{Results}
The reconstruction loss decreased monotonically from 4.44~nats to
\textbf{1.93~nats}; see Fig.~\ref{fig:t1}.
Five MIDI samples are exported to \texttt{outputs/generated\_midis/task1\_ae\_*}.

\begin{figure}[t]
\centering
\includegraphics[width=0.9\linewidth]{plots/task1_ae_loss.png}
\caption{Task 1 LSTM-AE reconstruction loss.}
\label{fig:t1}
\end{figure}

% ===========================================================================
\section{Task 2 --- $\beta$-Variational Autoencoder}
% ===========================================================================

\subsection{Formulation}
We extend Task~1 to the variational setting. The encoder outputs
$\mu(x), \log\sigma^2(x)$ and the decoder consumes $[z;\;g]$ where
$g\in\{0,1\}^{5}$ is a one-hot genre label:
\begin{align}
q_\phi(z\mid x) &= \mathcal{N}\!\big(\mu(x),\,\mathrm{diag}(\sigma^2(x))\big),\\
z &= \mu + \sigma \odot \epsilon,\quad \epsilon \sim \mathcal{N}(0,I).
\end{align}

\subsection{Loss}
\begin{equation}
\mathcal{L}_{VAE} = \mathcal{L}_{recon} + \beta\, D_{KL}\!\left(q_\phi(z\mid x)\,\Vert\, p(z)\right)
\label{eq:vae}
\end{equation}
with $p(z)=\mathcal{N}(0,I)$ and $\beta=0.05$.

\subsection{Training setup}
All five genres (282~windows), 30~epochs, batch~16, Adam with learning
rate $3\!\times\!10^{-3}$; 17.1~s total.

\subsection{Results}
Total loss fell from 4.39 to \textbf{2.64}, reconstruction from 4.30 to
2.53, while the KL term stabilised around 1.7--2.2~nats --- evidence that the
latent space is actively used and has not collapsed.
Fig.~\ref{fig:t2a} shows the three loss curves, and Fig.~\ref{fig:t2b}
visualises a linear interpolation between two $z$ vectors: the decoded
piano-rolls morph smoothly, confirming that nearby latent codes decode to
related patterns. Eight MIDI samples are exported
(\texttt{task2\_vae\_<genre>\_<i>.mid}).

\begin{figure}[t]
\centering
\includegraphics[width=0.95\linewidth]{plots/task2_vae_loss.png}
\caption{Task 2 VAE training: total, reconstruction and KL terms.}
\label{fig:t2a}
\end{figure}

\begin{figure}[t]
\centering
\includegraphics[width=0.95\linewidth]{plots/task2_latent_interp.png}
\caption{Latent interpolation: each row is the decoded piano-roll for
$z=(1-\alpha)z_1+\alpha z_2$ at five values of $\alpha$.}
\label{fig:t2b}
\end{figure}

% ===========================================================================
\section{Task 3 --- Transformer Decoder}
% ===========================================================================

\subsection{Architecture}
A GPT-style decoder-only Transformer with token + positional + genre
embeddings (all summed), two blocks of masked multi-head self-attention
(4~heads, model dimension $D=48$) and a linear output head. Each block
contains a residual-connection layer-norm pair and a feed-forward
sub-layer with hidden width $4D$.

\subsection{Loss and metric}
Standard next-token negative log-likelihood
\begin{equation}
\mathcal{L}_{TR}=-\frac{1}{T}\sum_{t=1}^{T}\log p_\theta(x_t\mid x_{<t}),
\end{equation}
with perplexity $\mathrm{PPL}=\exp(\mathcal{L}_{TR})$.

\subsection{Training setup}
All 282 windows, 10~epochs, batch~16, learning rate $2\!\times\!10^{-3}$;
53.0~s total.

\subsection{Results}
Cross-entropy dropped from 3.76 to \textbf{1.78} and perplexity from
\textbf{42.8 to 5.96} (Fig.~\ref{fig:t3}). Ten genre-conditioned
MIDI samples are exported; Fig.~\ref{fig:t3_rolls} shows one piano-roll
per genre.

\begin{figure}[t]
\centering
\includegraphics[width=0.95\linewidth]{plots/task3_tr_loss.png}
\caption{Task 3 Transformer: cross-entropy and perplexity.}
\label{fig:t3}
\end{figure}

\begin{figure}[t]
\centering
\includegraphics[width=0.98\linewidth]{plots/task3_pianorolls.png}
\caption{Piano-rolls of generated compositions for five genres
(Task 3 Transformer).}
\label{fig:t3_rolls}
\end{figure}

% ===========================================================================
\section{Task 4 --- RLHF Fine-Tuning}
% ===========================================================================

\subsection{Reward model}
A 10-listener preference survey of 20~clips is aggregated into a 3-feature
linear regression (in-scale ratio, rhythm variety, 1 minus repetition):
$r_{\mathrm{human}}(X) = w^\top f(X) + b$, fit by least squares.
The \emph{combined} reward also keeps a rule-based musicality prior
(scale fit, rhythm diversity, inverse repetition) with weight 0.6:
\begin{equation}
r(X) \;=\; 0.6\,r_{\mathrm{rule}}(X) \;+\; 0.4\,r_{\mathrm{human}}(X).
\end{equation}

\paragraph*{Data transparency} Because a real listening study with ten
human participants was not feasible within the project time-box, the
survey~CSV accompanying this report was produced by ten
\emph{simulated} listeners whose scores are noisy linear combinations of
the same three musicality features. The RLHF pipeline itself is unchanged:
the reward model $r_{\mathrm{human}}$ is fit from the CSV as it would be from
real scores. Replacing the CSV with genuine survey data only changes the
numerical reward; no code has to change.

\subsection{Algorithm}
REINFORCE with a batch-mean baseline~$b$ for variance reduction:
\begin{equation}
\nabla_\theta J(\theta) = \mathbb{E}\!\left[(r(X)-b)\sum_{t=1}^{T}\nabla_\theta\log p_\theta(x_t\mid x_{<t})\right].
\end{equation}

\subsection{Training setup}
20 RL steps, 4 sampled sequences per step, Adam lr $3\!\times\!10^{-4}$;
19.7~s total. The generator is initialised from the Task~3 checkpoint.

\subsection{Results}
Pre-RL combined reward averaged \textbf{0.557} across 6 samples; post-RL
reward averaged \textbf{0.656} across 10 samples, an increase of 18~\%
(Fig.~\ref{fig:t4}). Ten fine-tuned MIDI files are exported
(\texttt{task4\_rlhf\_<genre>\_<i>.mid}) plus five before-RLHF
controls (\texttt{task4\_before\_<genre>.mid}).

\begin{figure}[t]
\centering
\includegraphics[width=0.48\linewidth]{plots/task4_rlhf_reward.png}
\hfill
\includegraphics[width=0.48\linewidth]{plots/task4_before_after.png}
\caption{Left: mean reward across RL steps. Right: reward before and
after RLHF fine-tuning.}
\label{fig:t4}
\end{figure}

% ===========================================================================
\section{Baselines}
% ===========================================================================

Two non-neural baselines are implemented for comparison.

\subsubsection{Random-note generator}
Samples every token i.i.d. from the empirical unigram distribution of the
training corpus.

\subsubsection{First-order Markov chain}
Estimates transition matrix $P(x_t\mid x_{t-1})$ with add-one smoothing;
five MIDI samples were generated by ancestral sampling.

% ===========================================================================
\section{Evaluation}
% ===========================================================================

\subsection{Metrics}
Following the project rubric we report four quantitative measures and one
qualitative measure.
\begin{itemize}
  \item \textbf{Pitch histogram distance} --- $L_1$ distance between the
        12-bin pitch-class histograms of sample pairs (mean over all pairs).
        Higher values indicate that generated samples cover a wider part
        of pitch space.
  \item \textbf{Rhythm diversity} --- $\#\text{unique durations}\,/\,\#\text{total notes}$.
  \item \textbf{Repetition ratio} --- $1-(\#\text{unique }k\text{-grams}\,/\,\#k\text{-grams})$
        with $k\!=\!4$. Too-low values mean noisy output; too-high values
        mean boring loops.
  \item \textbf{Perplexity} --- only well-defined for Task~3 (an explicit
        autoregressive likelihood).
  \item \textbf{Human-listening score} --- the simulated 10-listener
        survey used to fit the RLHF reward model.
\end{itemize}

\subsection{Quantitative results}
Table~\ref{tab:results} summarises every model on the three objective
metrics; Fig.~\ref{fig:bars} visualises the same data.

\begin{table}[!t]
\centering
\caption{Quantitative comparison of baselines and task models.}
\label{tab:results}
\resizebox{\linewidth}{!}{
\begin{tabular}{l r r r l l}
\toprule
Model & Rhythm div. $\uparrow$ & Repetition $\downarrow/\uparrow^\ast$ & Pitch-hist. dist. & Perplexity & Genre control \\
\midrule
Random baseline & 0.036 & 0.000 & 0.476 & --- & none \\
Markov chain & 0.148 & 0.052 & 0.734 & --- & weak \\
Task 1 AE & 0.070 & 0.007 & 0.675 & --- & single \\
Task 2 VAE & 0.046 & 0.002 & 0.826 & --- & moderate \\
Task 3 Transformer & 0.217 & 0.151 & 1.070 & 5.96 & strong \\
Task 4 RLHF & 0.225 & 0.169 & 1.076 & --- & strong \\
\bottomrule
\end{tabular}
}
\end{table}

\begin{figure}[t]
\centering
\includegraphics[width=\linewidth]{plots/all_models_metrics.png}
\caption{Bar-chart comparison of the three objective metrics across all
six systems.}
\label{fig:bars}
\end{figure}

\subsection{Discussion}
\begin{itemize}
  \item The Random baseline has zero repetition because every output is
        effectively independent noise; yet its rhythm diversity
        (0.036) is far \emph{lower} than Transformer's (0.217),
        because the unigram distribution concentrates mass on frequent
        tokens.
  \item The Markov chain improves rhythm diversity but also introduces
        spurious repetitions and has no notion of genre structure beyond
        the marginal transitions.
  \item The LSTM-AE and VAE preserve a narrow pitch range because their
        loss averages over whole sequences --- they optimise fidelity,
        not variety. The VAE's slightly lower rhythm diversity is a
        consequence of its smoother latent prior.
  \item The Transformer is the first model to combine \emph{variety}
        (highest rhythm diversity) with \emph{structure}
        (non-trivial repetition ratio), reflected in its low perplexity
        of~5.96.
  \item RLHF fine-tuning shifts Transformer outputs toward higher-reward
        patterns, measurable as +0.10 in combined reward and a small
        additional gain in rhythm diversity, at the cost of slightly
        more repetition --- a well-known effect of reward over-optimisation.
\end{itemize}

% ===========================================================================
\section{Limitations and Future Work}
% ===========================================================================

Due to compute constraints we limit $T\!=\!64$ 16\textsuperscript{th} notes
(four bars) and train for tens of epochs; scaling to MAESTRO-sized
corpora would use the same code-path but warrant bigger models and GPU
training (e.g. a PyTorch drop-in using the identical tokeniser).
The monophonic reduction loses harmony information; a polyphonic
piano-roll output head is an immediate extension. The simulated listener
survey should be replaced by a real one before any claim is made about
human preference; the code path is deliberately designed so that only the
CSV file changes.

% ===========================================================================
\section{Conclusion}
% ===========================================================================

We presented a single, reproducible NumPy+Autograd codebase that trains
four progressively more sophisticated unsupervised music generators --- an
LSTM autoencoder, a $\beta$-VAE, a Transformer decoder and an RLHF
fine-tuned policy --- on a shared multi-genre MIDI tokeniser. The
Transformer attains a per-token perplexity of~5.96 after 53~seconds of
CPU training, and RLHF raises the combined musicality reward by~18\%.
All 48~generated MIDI samples, loss curves and evaluation tables are
packaged with the source tree.

\begin{thebibliography}{99}
\bibitem{maestro} C. Hawthorne, et al. ``Enabling Factorized Piano Music Modeling and Generation with the MAESTRO Dataset,'' \emph{ICLR}, 2019.
\bibitem{lakh} C. Raffel. ``Learning-Based Methods for Comparing Sequences, with Applications to Audio-to-MIDI Alignment and Matching,'' Ph.D.\ dissertation, Columbia Univ., 2016.
\bibitem{groove} J. Gillick, et al. ``Learning to Groove with Inverse Sequence Transformations,'' \emph{ICML}, 2019.
\bibitem{lstm} S. Hochreiter and J. Schmidhuber. ``Long Short-Term Memory,'' \emph{Neural Computation}, vol.~9, no.~8, 1997.
\bibitem{vae} D. P. Kingma and M. Welling. ``Auto-Encoding Variational Bayes,'' \emph{ICLR}, 2014.
\bibitem{transformer} A. Vaswani, et al. ``Attention Is All You Need,'' \emph{NeurIPS}, 2017.
\bibitem{rlhf} P. Christiano, et al. ``Deep Reinforcement Learning from Human Preferences,'' \emph{NeurIPS}, 2017.
\bibitem{musegan} H.-W. Dong, et al. ``MuseGAN: Multi-track Sequential Generative Adversarial Networks for Symbolic Music Generation and Accompaniment,'' \emph{AAAI}, 2018.
\end{thebibliography}

\end{document}